\chapter{文献综述}
\section{已实现波动率预测研究}
\subsection{低频条件波动与高频已实现测度}
GARCH 模型通过收益创新和滞后条件方差描述波动动态，为低频收益的条件异方差建模提供了基本工具\cite{bollerslev1986}。高频已实现测度则从日内价格路径出发，将多个高频收益的平方累计为日度风险代理。两者并非互相排斥：条件波动模型刻画日收益分布，已实现测度用于构造和预测日内变化；本文分别将其用于风险边与预测目标。

在适当的价格过程及观测条件下，已实现方差与二次变差之间存在理论联系\cite{andersen2003}。实际数据仍涉及有限采样、价格噪声、停牌与不规则时段等问题。因此，高频数据行数更多并不自动意味着日度测度更准确；一个明确定义、来源一致的五分钟价格路径，通常比将不同时段和来源的记录机械相加更适合可复核比较。本文选择五分钟作为起始采样尺度，并以替代采样频率检验结果对测度的敏感性，不将其预设为所有数据场景的最优选择。

\subsection{HAR、测量误差与自身信息扩展}
HAR 将不同时间尺度的滞后已实现波动率放入简约回归，以近似表示波动率较长的持续性\cite{corsi2009}。其价值不仅体现在基准预测能力，还体现在结构清楚：任何新增公共因子、图聚合项或非线性分支，都可以与自身历史形成直接对照。

Bollerslev、Patton 和 Quaedvlieg 将已实现测度的误差大小纳入预测系数调整，提出 HARQ 等模型\cite{bpq2016}。本文采用其日度四次变差交互项的简约版本作为计量基线。HARQ 中的 Q 指四次变差相关的误差调整，并非 QLIKE 训练的缩写；后者在本文单独标注为训练损失。

对于自身信息，半方差或收益符号能够提供不同于总方差的描述。本文设置相应的自身特征扩展对照，是为了避免把本来可以由目标股票自身状态解释的改善全部归因于外部风险边。本研究不将所有自身特征扩展同时堆入主模型，而按照信息匹配原则在明确的比较组中使用。

\subsection{预测损失与模型比较}
方差通常通过代理量评价，损失函数对比较结论具有重要影响。Patton 讨论了使用不完美波动率代理时的预测比较条件\cite{patton2011}。这一研究为考察不同损失提供了依据，但不意味着任意来源、任意处理过的 RV 都自动满足相关条件。

在本文中，MSE 强调原始方差单位下的平方偏差，QLIKE 依赖实际与预测的比值。两者既是可能的训练准则，也是需要明确区分的评价准则。模型使用 QLIKE 训练之后，仍需与使用同一准则训练的无图模型比较。由此获得的差值，才更接近网络表征的增量，而不是训练方式改变的混合结果。

\section{波动率溢出与动态 Copula-CoVaR 研究}
\subsection{条件风险与依赖结构}
多资产风险建模需要区分单个资产的边际分布与资产之间的依赖关系。动态相关模型为时间变化的关联提供简约参数化\cite{engle2002}；Copula 表示则将边际分布与联合依赖分开，使不同边际下的条件风险计算可以在统一的概率坐标中进行。

Adrian 和 Brunnermeier 提出的 CoVaR 从某一主体处于特定状态的条件出发，考察另一主体或系统的风险分位数，增量形式则比较困境与基准状态\cite{adrian2016}。本文将这一思路用于股票对之间的条件风险关系，但研究对象和目标不同于原有系统性风险度量：股票对风险边在这里是未来已实现方差的预测输入，而不是对宏观系统重要性的直接认定。

\subsection{条件定义与方向语义}
来源收益恰好处于一个分位数位置，与来源收益落入该分位数以下的整个尾部事件，是两个不同条件。Mainik 和 Schaanning 讨论了不同条件设定下 CoVaR 对依赖关系的一致性问题\cite{mainik2012}。本文因此将点条件与事件条件分开定义、保存和比较，不在程序中只用一个 CoVaR 名称掩盖二者差异。

条件方向也不同于时间领先或因果方向。$i\leftarrow j$ 表示在$j$的给定分位状态下计算$i$的条件分位数；即使同时得到两个方向数值不同，也还需要分析这种差异是否由边际尺度造成。本文所称“风险溢出网络”沿用开题中的经济解释，但对可识别内容作出上述限制。

\subsection{尾部对称性与非对称扩展}
标准 t-Copula 具有上下尾对称的结构\cite{demarta2005}。在对称边际和镜像分位条件下，仅将分位参数从下尾改到上尾并不会产生独立的第二种风险机制。本文把这一性质写入模型定义与前置测试，以防将数值误差或重复变量解释成市场非对称证据。

允许上下尾不同的 Copula 可以用于进一步研究，例如 Patton 对非对称依赖的建模\cite{patton2006}。本文将此类模型作为明确命名的扩展，而不是在原型效果不理想后无记录地替换。上下尾不对称本身也不意味着来源与目标之间的非交换性，两者需分别检验。

\subsection{风险汇总与已实现波动率预测}
与本文最直接相邻的工作是 Wang、Yuan 和 Qu 的极端市场状态切换模型。公开摘要与引言显示，该研究基于动态 Copula-CoVaR 构建上、下行风险输出和接收指标，并将其与 HAR 相结合\cite{wang2026}。这一证据已经说明，CoVaR 与波动率预测的结合并非本文首次提出。

本文的区别在于进一步询问：当节点输入、输出的风险总量已被模型使用后，具体邻居与边特征之间的对应关系是否仍有价值。由于目前未取得该工作的完整方法公式与实现，本文先构造定义透明的风险汇总基线，不将其命名为该文的严格复现；取得全文后再通过公式对应表补充复现版本。

\section{图神经网络与金融图预测}
\subsection{图结构与边输入}
图模型把股票作为节点，通过邻居聚合描述横截面联系。GLASSO 在精度矩阵上进行稀疏估计，是统计构图的一种方式\cite{friedman2008}。其所反映的是特定模型假设下的条件关联结构，不应被直接解释为因果关系。本文保留其作为竞争构图方式，而不是先验否定其风险描述能力。

消息传递框架允许模型根据节点和边的信息构造消息并聚合到目标节点\cite{gilmer2017}。因此，边特征机制本身不是本文发明的通用图学习算子。本文的方法贡献在于为这一任务选择明确的风险边表示，并在统一实验中检验其是否比标量权重或风险汇总更有用。

\subsection{GNNHAR 与训练准则}
Zhang 等将图上的非线性溢出纳入多资产已实现波动率预测，并同时研究邻居范围和训练准则\cite{zhang2025}。这一工作提示本文应把信息集合、函数形式和训练损失分别控制，而不把它们的联合改变压缩成单一的“GNN 改进”。

本文设置两组相关对照：一组尽可能对应原文的图与模型设定，另一组在统一输出约束、节点变量和训练预算下替换构图及边表示。前者用于说明与既有文献的联系，后者用于更精确地识别本项目内部的增量来源。预印本与最终发表版的实现差异亦需在复现说明中记录。

\section{文献评述与本文定位}
\begin{table}[htbp]\centering\small
\caption{相关研究与本文研究问题的对应关系}
\begin{tabularx}{\textwidth}{p{0.2\textwidth}YY}
\toprule 研究线索 & 已提供的基础 & 本文进一步检验的内容\\\midrule
HAR 与 HARQ & 多尺度自身历史、测量误差调整 & 新增横截面信息能否超过强自身基线\\
GNNHAR & 图上的非线性聚合及损失比较 & 风险图与多维边是否有独立增量\\
CoVaR 风险汇总预测 & 风险输出/接收及状态信息 & 同风险汇总下保留具体关系的价值\\
Copula 风险度量 & 边际与依赖分离、条件分位数 & 风险语义、尺度与模型识别边界\\
边特征消息传递 & 关系依赖的消息函数 & 对这一金融任务的有限样本适用性\\\bottomrule
\end{tabularx}
\end{table}
本文不以“已有方法都忽略尾部风险”或“相关图无法识别任何风险”作为出发点。更具体的定位是：在既有预测与风险度量基础上，建立从风险汇总到逐边条件化的可检验联系。若新增表示能够稳定提高预测能力，说明这种关系信息在所研究样本与协议下值得保留；若风险汇总已足够，则说明更简约的信息压缩在相应条件下具有竞争力。两类结果都应通过相同实验标准判断。
